%============================================================
% CH1 · PATTERN 3 — GRIDS AS GRAPHS: THEORY
%============================================================
% Local helpers of Pattern 3 (TikZ styles g1p3..., macros \gIpIII...):
%   \gIpIIIgrid{row,row,...}  draws a grid inside a tikzpicture; one character per cell:
%       o open   w wall/water   A..F component colours with the letter, a..f same colours without it
%       s start cell   x current cell   q on the stack / in the queue   k finished
%   cell (r,c) is centred at (\gIpIIIs*c, -\gIpIIIs*r);  \gIpIIIat{r}{c}{text} writes into a cell
\tikzset{
    g1p3o/.style={gcell},
    g1p3w/.style={gwall},
    g1p3a/.style={gcell, fill=primaryblue!22},
    g1p3b/.style={gcell, fill=successgreen!25},
    g1p3c/.style={gcell, fill=accentorange!35},
    g1p3d/.style={gcell, fill=softpurple!25},
    g1p3e/.style={gcell, fill=recursionred!22},
    g1p3f/.style={gcell, fill=insightborder!30},
    g1p3A/.style={g1p3a, label={[font=\scriptsize\bfseries, text=darkgray]center:A}},
    g1p3B/.style={g1p3b, label={[font=\scriptsize\bfseries, text=darkgray]center:B}},
    g1p3C/.style={g1p3c, label={[font=\scriptsize\bfseries, text=darkgray]center:C}},
    g1p3D/.style={g1p3d, label={[font=\scriptsize\bfseries, text=darkgray]center:D}},
    g1p3E/.style={g1p3e, label={[font=\scriptsize\bfseries, text=darkgray]center:E}},
    g1p3F/.style={g1p3f, label={[font=\scriptsize\bfseries, text=darkgray]center:F}},
    g1p3s/.style={gcell, fill=basefill, draw=successgreen, line width=1.2pt},
    g1p3x/.style={gcell, fill=currentfill, draw=accentorange, line width=1.2pt},
    g1p3q/.style={gcell, fill=hitfill},
    g1p3k/.style={gcell, fill=computedfill},
    g1p3lbl/.style={font=\scriptsize\bfseries, text=darkgray, inner sep=0pt},
    g1p3var/.style={rectangle, rounded corners=2pt, draw=patternborder, fill=patternbg,
                    font=\scriptsize, align=center, text width=3.25cm, minimum height=1.25cm, inner sep=3pt}
}
\newcount\gIpIIIc
\def\gIpIIIs{0.7}
\def\gIpIIIscan#1{\ifx#1\relax\else
    \node[g1p3#1, minimum size=\gIpIIIs cm] at (\gIpIIIs*\the\gIpIIIc, -\gIpIIIs*\gIpIIIr) {};%
    \global\advance\gIpIIIc by 1 \expandafter\gIpIIIscan\fi}
\newcommand{\gIpIIIgrid}[1]{%
    \foreach \gIpIIIrowtxt [count=\gIpIIIr from 0] in {#1}{%
        \global\gIpIIIc=0 \expandafter\gIpIIIscan\gIpIIIrowtxt\relax}}
\newcommand{\gIpIIIat}[3]{\node[g1p3lbl] at (\gIpIIIs*#2, -\gIpIIIs*#1) {#3};}

\subsection{Pattern Theory}
\label{ssec:ch1-p3-theory}

\begin{patternbox}{Pattern 3 at a Glance}
A grid \emph{is} a graph: every open \textbf{cell is a vertex}, and two open cells that \textbf{share a side are joined by an edge}. We never store this graph --- the neighbours of $(r,c)$ are computed on the fly. A ``room'', ``island'', ``lake'' or ``region'' is a \textbf{connected component}, found by a \textbf{flood fill} (a DFS or BFS that spreads from one cell to its open neighbours).

\begin{center}
\begin{tikzpicture}
    \begin{scope}
        \def\gIpIIIs{0.75}
        \gIpIIIgrid{ooo,oxo,ooo}
        \draw[flow] (0.75,-0.75) -- (0.75,-0.1);
        \draw[flow] (0.75,-0.75) -- (0.75,-1.4);
        \draw[flow] (0.75,-0.75) -- (0.1,-0.75);
        \draw[flow] (0.75,-0.75) -- (1.4,-0.75);
        \node[font=\scriptsize] at (0.75,0.65) {$(r{-}1,c)$};
        \node[font=\scriptsize] at (0.75,-2.15) {$(r{+}1,c)$};
        \node[font=\scriptsize, anchor=east] at (-0.45,-0.75) {$(r,c{-}1)$};
        \node[font=\scriptsize, anchor=west] at (1.95,-0.75) {$(r,c{+}1)$};
    \end{scope}
    \begin{scope}[xshift=6.3cm, yshift=0.3cm]
        \def\gIpIIIs{0.6}
        \gIpIIIgrid{AAwwB,wAwBB,CwwwB,CCwDw}
        \node[font=\scriptsize, align=center] at (1.2,-2.75) {4 components: sizes 3, 4, 3, 1};
    \end{scope}
\end{tikzpicture}
\end{center}

\begin{recurrencebox}
\small
\textbf{for} every cell $(r,c)$ in row-major order: \textbf{if} it is open and unvisited \textbf{then}\\
\hspace*{1.5em}\texttt{components++}; mark $(r,c)$; push it on a stack (or queue)\\
\hspace*{1.5em}\textbf{while} the stack is not empty: pop $(x,y)$; \texttt{size++};\\
\hspace*{3em}\textbf{for} each of the 4 neighbours $(x{+}dr_k,\,y{+}dc_k)$ that is \emph{inside}, \emph{open} and \emph{unvisited}: mark it, push it.\\[2pt]
\hspace*{\fill}$\bigO{n\cdot m}$ time, $\bigO{n\cdot m}$ memory --- every cell is marked once and looks at 4 neighbours.
\end{recurrencebox}
\end{patternbox}

%------------------------------------------------------------
\subsubsection*{From Pattern 2 to Pattern 3}

In Pattern 2 the graph was \emph{handed to us}: $n$ vertices and a list of edges, stored in an adjacency list, and we asked ``can $s$ reach $t$?''. In this pattern the input is a \textbf{picture} --- a map of $n$ rows and $m$ columns of characters or numbers --- and nobody mentions vertices or edges. The graph is still there; we only have to \emph{see} it:

\begin{center}
\renewcommand{\arraystretch}{1.25}
\begin{tabular}{@{}ll@{}}
\toprule
\textbf{Picture word} & \textbf{Graph word} \\
\midrule
open cell (floor `\texttt{.}', land \texttt{1}, a pixel of colour $x$) & vertex \\
two open cells sharing a side & edge (undirected) \\
``you can walk from $A$ to $B$'' & $A$ and $B$ are in the same connected component \\
room / island / lake / region / blob & connected component \\
area of an island & number of vertices of a component \\
number of rooms & number of components \\
\bottomrule
\end{tabular}
\end{center}

\begin{diagrambox}
\begin{tikzpicture}
    \begin{scope}
        \gIpIIIgrid{oowo,owoo,oooo}
        \foreach \c in {0,1,2,3} \node[font=\scriptsize, text=darkgray] at (0.7*\c, 0.55) {$c{=}\c$};
        \foreach \r in {0,1,2} \node[font=\scriptsize, text=darkgray, anchor=east] at (-0.45, -0.7*\r) {$r{=}\r$};
        \node[font=\small\bfseries] at (1.05,-2.25) {the map ($3\times4$, two walls)};
    \end{scope}
    \draw[-{Stealth[length=3mm]}, line width=1.4pt, primaryblue] (3.0,-0.7) -- node[above, font=\scriptsize]{see it as} (4.4,-0.7);
    \begin{scope}[xshift=5.2cm]
        \foreach \r/\c in {0/0,0/1,0/3,1/0,1/2,1/3,2/0,2/1,2/2,2/3}
            \node[gv, minimum size=6.5mm, font=\scriptsize] (v\r\c) at (1.25*\c, -1.0*\r) {\r,\c};
        \foreach \a/\b in {00/01, 00/10, 03/13, 10/20, 12/13, 12/22, 13/23, 20/21, 21/22, 22/23}
            \draw[ge] (v\a) -- (v\b);
        \node[font=\small\bfseries] at (1.9,-2.75) {the graph: 10 vertices, 10 edges};
    \end{scope}
\end{tikzpicture}
\figcap{Each open cell becomes a vertex named by its coordinates; walls simply have no vertex. An edge joins every pair of open cells that share a side. Cells $(0,1)$ and $(1,0)$ touch only at a corner: no edge.}
\end{diagrambox}

\textbf{How big is this graph?} At most $n\cdot m$ vertices, and every vertex has at most 4 neighbours, so $|E| \le 4nm/2 = 2nm$ (a fully open grid has exactly $n(m-1) + m(n-1)$ edges). For the usual limits $n, m \le 1000$ that is $10^6$ vertices and $2\cdot10^6$ edges --- well within an $\bigO{|V|+|E|}$ algorithm, but \emph{building} an adjacency list for it would waste time and tens of megabytes on information we can recompute instantly. So we keep the grid itself as the graph: it is an \textbf{implicit graph}, exactly the kind described in Lesson~\ref{ssec:ch1-modelling}.

%------------------------------------------------------------
\subsubsection*{Neighbours with two small arrays}

Instead of four copies of the same \texttt{if}, store the four moves as \emph{direction vectors} and loop over them:

\begin{diagrambox}
\begin{tikzpicture}
    \begin{scope}
        \gIpIIIgrid{ooo,oxo,ooo}
        \foreach \x/\y/\t in {0.7/0/{$-1,0$}, 0.7/-1.4/{$+1,0$}, 0/-0.7/{$0,-1$}, 1.4/-0.7/{$0,+1$}}
            \node[font=\tiny\bfseries, text=deepblue] at (\x,\y) {\t};
        \node[font=\scriptsize] at (0.7,-0.7) {$r,c$};
        \node[font=\small\bfseries, align=center] at (0.7,-2.15) {4-neighbourhood};
        \node[font=\scriptsize, align=center] at (0.7,-2.8) {\texttt{dr = \{-1, 1, 0, 0\}}\\\texttt{dc = \{ 0, 0,-1, 1\}}};
    \end{scope}
    \begin{scope}[xshift=4.6cm]
        \gIpIIIgrid{qqq,qxq,qqq}
        \foreach \x/\y/\t in {0/0/{$-1,-1$}, 0.7/0/{$-1,0$}, 1.4/0/{$-1,+1$}, 0/-0.7/{$0,-1$}, 1.4/-0.7/{$0,+1$}, 0/-1.4/{$+1,-1$}, 0.7/-1.4/{$+1,0$}, 1.4/-1.4/{$+1,+1$}}
            \node[font=\tiny\bfseries, text=deepblue] at (\x,\y) {\t};
        \node[font=\scriptsize] at (0.7,-0.7) {$r,c$};
        \node[font=\small\bfseries, align=center] at (0.7,-2.15) {8-neighbourhood};
        \node[font=\scriptsize, align=center] at (0.7,-2.8) {add the 4 diagonals\\(``king moves'')};
    \end{scope}
    \begin{scope}[xshift=9.2cm]
        \gIpIIIgrid{ooooo,ooooo,ooxoo,ooooo,ooooo}
        \foreach \r/\c in {0/1,0/3,1/0,1/4,3/0,3/4,4/1,4/3} \node[g1p3q, minimum size=0.7cm] at (0.7*\c,-0.7*\r) {};
        \node[font=\scriptsize] at (1.4,-1.4) {$r,c$};
        \node[font=\small\bfseries, align=center] at (1.4,-3.55) {knight moves};
        \node[font=\scriptsize, align=center] at (1.4,-4.1) {$(\pm1,\pm2),\ (\pm2,\pm1)$};
    \end{scope}
\end{tikzpicture}
\figcap{The neighbourhood is part of the \emph{model}, not of the algorithm. Read the statement: ``share a side'' / ``up, down, left, right'' = 4 neighbours; ``touch'' / ``including diagonally'' = 8; a chess piece = its own move list. The flood fill code stays identical --- only the arrays change.}
\end{diagrambox}

The neighbour test has three parts, and their \textbf{order matters}: first \emph{inside the grid} (otherwise \texttt{g[nr][nc]} reads outside the array), then \emph{open}, then \emph{not yet visited}.

\subsubsection*{Naming a cell with one integer}

Pairs $(r,c)$ are readable, but sometimes one integer per vertex is handier (a \texttt{vector<int>} stack, a DSU, a \texttt{visited} array of size $nm$). Number the cells row by row:
\[
\text{id} = r\cdot m + c, \qquad r = \text{id} \;/\; m, \qquad c = \text{id} \bmod m \qquad (0 \le \text{id} < nm).
\]
\begin{diagrambox}
\begin{tikzpicture}
    \gIpIIIgrid{oooo,oooo,oooo}
    \foreach \r in {0,1,2} \foreach \c in {0,1,2,3} {
        \pgfmathtruncatemacro{\id}{\r*4+\c}
        \gIpIIIat{\r}{\c}{\id}
    }
    \foreach \c in {0,1,2,3} \node[font=\scriptsize, text=darkgray] at (0.7*\c, 0.55) {$c{=}\c$};
    \foreach \r in {0,1,2} \node[font=\scriptsize, text=darkgray, anchor=east] at (-0.45, -0.7*\r) {$r{=}\r$};
    \node[note, anchor=west, text width=7.2cm, align=left] at (3.0,-0.7) {$n=3$ rows, $m=4$ columns. Cell $(2,1)$ has id $2\cdot4+1 = 9$; id $6$ is $(6/4,\ 6 \bmod 4) = (1,2)$. Multiply by the number of \textbf{columns} $m$, not by $n$ --- with $n \ne m$ the wrong factor silently mixes cells up.};
\end{tikzpicture}
\end{diagrambox}

%------------------------------------------------------------
\subsubsection*{When to recognise it}

\begin{itemize}[leftmargin=*]
    \item the input is a \textbf{2D map}: characters (\texttt{.} / \texttt{\#}), 0/1 values, colours, heights, depths;
    \item movement is between \textbf{adjacent cells} (4 or 8 directions, or a fixed move set);
    \item the question is about \textbf{regions}: count them (rooms, islands), measure them (area, volume, perimeter), recolour one (flood fill), or decide whether two cells are in the same region.
\end{itemize}
\textbf{Not this pattern (yet):}
\begin{itemize}[leftmargin=*]
    \item ``\textbf{minimum number of steps}'' from one cell to another: the order of exploration matters, use BFS with distances --- Chapter 3 (BFS);
    \item cells with \textbf{different costs} to enter: Dijkstra or 0-1 BFS, later chapters;
    \item the grid \textbf{changes} over time (walls added or removed between queries): Chapter 6 (DSU);
    \item the cell alone is not enough to describe where you are (you carry keys, you have used $k$ wall-breaks, you face a direction): the vertex becomes (cell, extra information) --- a state graph, Pattern 4.
\end{itemize}

%------------------------------------------------------------
\subsubsection*{The algorithm: scan and fill}

\begin{enumerate}[leftmargin=*]
    \item Create \texttt{visited} (or reuse the grid: overwrite a visited open cell with a wall character).
    \item Scan the cells in row-major order. An open, unvisited cell is the \textbf{first cell of a new component}: count it and start a flood fill there.
    \item The flood fill marks the start cell and puts it on a stack. While the stack is not empty, pop a cell and, for each of its neighbours that is inside, open and unvisited, \textbf{mark it immediately} and push it.
    \item When the stack is empty, the whole component is marked; the scan continues and skips every marked cell.
\end{enumerate}

\begin{diagrambox}
\begin{tikzpicture}
    \begin{scope}
        \gIpIIIgrid{sowwo,wowoo,owwwo,oowow}
        \draw[flow, dashed] (-0.35,0.5) -- (3.15,0.5);
        \node[font=\scriptsize, text=primaryblue] at (1.4,0.75) {scan row by row};
        \node[font=\scriptsize\bfseries, align=center] at (1.4,-3.0) {scan reaches $(0,0)$:\\open, unvisited $\Rightarrow$ component A};
    \end{scope}
    \begin{scope}[xshift=4.6cm]
        \gIpIIIgrid{AAwws,wAwoo,owwwo,oowow}
        \node[font=\scriptsize\bfseries, align=center] at (1.4,-3.0) {A is filled (3 cells);\\next new start: $(0,4)$};
    \end{scope}
    \begin{scope}[xshift=9.2cm]
        \gIpIIIgrid{AAwwB,wAwBB,CwwwB,CCwDw}
        \foreach \r/\c in {0/0,0/4,2/0,3/3} \draw[successgreen, line width=1.3pt] (0.7*\c-0.35,-0.7*\r-0.35) rectangle (0.7*\c+0.35,-0.7*\r+0.35);
        \node[font=\scriptsize\bfseries, align=center] at (1.4,-3.0) {done: 4 starts (green frames)\\= 4 components, sizes 3, 4, 3, 1};
    \end{scope}
\end{tikzpicture}
\figcap{The scan finds the first cell of each component in row-major order; the fill then swallows the whole component, so the scan never starts twice inside it. The template below prints ``4 components, sizes: 3 4 3 1'' for this grid.}
\end{diagrambox}

\begin{invariantbox}[Why scan-and-fill counts components exactly]
\textbf{Claim 1 (a fill marks exactly one component).} A fill started at $s$ marks precisely the open cells reachable from $s$.
\emph{Only reachable cells:} a cell is marked only when it is a neighbour of a marked cell, so by induction every marked cell has a path from $s$.
\emph{Every reachable cell:} suppose a reachable cell $v$ stays unmarked; on a path $s = u_0, u_1, \dots, u_k = v$ take the first unmarked $u_{i+1}$. Its predecessor $u_i$ was marked, hence pushed, hence popped, and when it was popped all its open unvisited neighbours --- $u_{i+1}$ among them --- were marked. Contradiction.

\textbf{Claim 2 (one start per component).} A fill starts only at an unmarked cell; by Claim~1 that cell's component is still completely unmarked (fills mark whole components), and after the fill it is completely marked, so no later start lies in it. Every component gets a start because the scan visits all cells. Hence \#starts $=$ \#components.

\textbf{Cost.} Each cell is marked at most once and therefore pushed and popped at most once; each pop inspects 4 neighbours. Total $\bigO{nm}$ time and $\bigO{nm}$ memory for \texttt{visited} and the stack. \hfill$\blacksquare$
\end{invariantbox}

\subsubsection*{Stack or queue? Recursion?}

The fill works with \emph{any} container: the set of marked cells at the end is the same, only the \emph{order} differs.

\begin{diagrambox}
\begin{tikzpicture}
    \begin{scope}
        \gIpIIIgrid{kkkk,kskk,kkkk,kkkk}
        \foreach \r/\c/\t in {0/0/14,0/1/16,0/2/12,0/3/10, 1/0/13,1/1/1,1/2/2,1/3/3, 2/0/9,2/1/15,2/2/11,2/3/4, 3/0/8,3/1/7,3/2/6,3/3/5}
            \gIpIIIat{\r}{\c}{\t}
        \node[font=\small\bfseries, align=center] at (1.05,-3.2) {stack (DFS order)\\[-1pt]{\scriptsize\mdseries runs away along a long path}};
    \end{scope}
    \begin{scope}[xshift=4.6cm]
        \gIpIIIgrid{kkkk,kskk,kkkk,kkkk}
        \foreach \r/\c/\t in {0/0/6,0/1/2,0/2/7,0/3/12, 1/0/4,1/1/1,1/2/5,1/3/11, 2/0/9,2/1/3,2/2/10,2/3/15, 3/0/13,3/1/8,3/2/14,3/3/16}
            \gIpIIIat{\r}{\c}{\t}
        \node[font=\small\bfseries, align=center] at (1.05,-3.2) {queue (BFS order)\\[-1pt]{\scriptsize\mdseries grows in rings around the start}};
    \end{scope}
    \begin{scope}[xshift=9.2cm]
        \gIpIIIgrid{skkk,kkkk,kkkk,kkkk}
        \draw[gpath, opacity=0.5, -{Stealth[length=2.5mm]}] (0,0) -- (0,-2.1) -- (0.7,-2.1) -- (0.7,0) -- (1.4,0) -- (1.4,-2.1) -- (2.1,-2.1) -- (2.1,0);
        \foreach \r/\c/\t in {0/0/1,0/1/8,0/2/9,0/3/16, 1/0/2,1/1/7,1/2/10,1/3/15, 2/0/3,2/1/6,2/2/11,2/3/14, 3/0/4,3/1/5,3/2/12,3/3/13}
            \gIpIIIat{\r}{\c}{\t}
        \node[font=\small\bfseries, align=center] at (1.05,-3.2) {recursive DFS from $(0,0)$\\[-1pt]{\scriptsize\mdseries depth 16 = every cell!}};
    \end{scope}
\end{tikzpicture}
\figcap{Pop order on an open $4\times4$ grid (neighbour order up, down, left, right; green = start). Left and middle start at $(1,1)$. Right: the recursive DFS from $(0,0)$ snakes through all 16 cells before returning even once, so the recursion depth equals the number of cells.}
\end{diagrambox}

\begin{warningbox}[Recursion depth on grids]
The picture on the right is not a worst case, it is the \emph{normal} case: on an open $1000\times1000$ room the recursive fill reaches depth $10^6$ (we measured it: the program crashes with the default stack, while depth $10^4$ on a $100\times100$ grid is fine). Rule of thumb: recursion is safe for $nm$ up to about $10^5$ (LeetCode limits like $50\times50$ or $300\times300$ are fine); for $nm$ up to $10^6$ (CSES, Codeforces) use the \textbf{explicit stack or the queue}. The short recursive version is still worth knowing: it is the fastest to type when it is safe.
\end{warningbox}

%------------------------------------------------------------
\subsubsection*{Variations map}

\begin{diagrambox}
\begin{tikzpicture}
    \node[g1p3var, fill=answerfill, draw=accentorange, text width=4.6cm, font=\scriptsize\bfseries] (root) at (0,0)
        {Scan + flood fill\\cells = vertices, sides = edges};
    \node[g1p3var] (v1) at (-5.4,-2) {\textbf{count components}\\rooms, islands\\(Problem 1.3.1)};
    \node[g1p3var] (v2) at (-1.8,-2) {\textbf{size / sum per component}\\largest island, lake volume\\(Problem 1.3.2, E3.3)};
    \node[g1p3var] (v3) at (1.8,-2)  {\textbf{fill from one cell}\\recolour a region\\(Flood Fill, E3.2)};
    \node[g1p3var] (v4) at (5.4,-2)  {\textbf{boundary of a region}\\perimeter = sides facing\\water (E3.1)};
    \node[g1p3var] (v5) at (-5.4,-4.1) {\textbf{a property of a whole component}\\``every cell is also land in grid 1''\\(E3.4)};
    \node[g1p3var] (v6) at (-1.8,-4.1) {\textbf{label + sizes, then combine}\\flip one cell to join islands\\(E3.5)};
    \node[g1p3var] (v7) at (1.8,-4.1)  {\textbf{fill from the border}\\regions that cannot escape\\(Chapter 4)};
    \node[g1p3var] (v8) at (5.4,-4.1)  {\textbf{shortest path in a maze}\\BFS with distances\\(Chapter 3)};
    \node[draw=patternborder, dashed, rounded corners=4pt, fit=(v1)(v4)(v5)(v8), inner sep=5pt] (grp) {};
    \draw[flow] (root.south) -- (grp.north);
\end{tikzpicture}
\end{diagrambox}

\textbf{Twists and how to handle them}
\begin{itemize}[leftmargin=*]
    \item \textbf{8-connectivity} (``cells touching at a corner are connected''): use the 8-entry direction arrays; nothing else changes.
    \item \textbf{Same colour instead of open/blocked} (images, height maps): ``open'' becomes ``has the same value as the start cell''; store the start value before you overwrite anything.
    \item \textbf{Wrap-around} (a torus: leaving on the right enters on the left): replace the inside test by \texttt{nr = (r + dr[d] + n) \% n} and \texttt{nc = (c + dc[d] + m) \% m}.
    \item \textbf{Two grids at once} (sub-islands, comparing pictures): fill in one grid while checking the other.
    \item \textbf{Per-component information used later} (size of the island next to a water cell): first label every cell with its component id and store the sizes, then answer each question in $\bigO{1}$.
    \item \textbf{Many test cases with a total size limit}: allocate \texttt{visited} per test with the real $n\times m$, never a global $1000\times1000$ array cleared every time.
\end{itemize}

%------------------------------------------------------------
\subsubsection*{Code templates}

The toolkit below contains everything this pattern needs: the inside test, the direction arrays, the recursive fill (small grids only), the iterative fill with an explicit stack, the BFS fill, and the scan that labels all components.

\begin{lstlisting}[style=cpp]
#include <bits/stdc++.h>
using namespace std;

int n, m;                                    // rows, columns
vector<string> g;                            // '.' = open cell, '#' = blocked
const int dr[4] = {-1, 1, 0, 0};             // up, down, left, right
const int dc[4] = {0, 0, -1, 1};

bool inside(int r, int c) { return 0 <= r && r < n && 0 <= c && c < m; }
bool open(int r, int c) { return inside(r, c) && g[r][c] == '.'; }

// (1) recursive DFS: fine when the grid is small (depth can reach n*m!)
void dfs(int r, int c, int id, vector<vector<int>>& comp) {
    comp[r][c] = id;
    for (int d = 0; d < 4; d++) {
        int nr = r + dr[d], nc = c + dc[d];
        if (open(nr, nc) && comp[nr][nc] == -1) dfs(nr, nc, id, comp);
    }
}

// (2) iterative DFS with an explicit stack: safe for 10^6 cells; returns the size
int fillStack(int sr, int sc, int id, vector<vector<int>>& comp) {
    vector<pair<int,int>> st;
    comp[sr][sc] = id;                       // mark when PUSHED
    st.push_back(make_pair(sr, sc));
    int size = 0;
    while (!st.empty()) {
        int r = st.back().first, c = st.back().second;
        st.pop_back();
        size++;
        for (int d = 0; d < 4; d++) {
            int nr = r + dr[d], nc = c + dc[d];
            if (!open(nr, nc) || comp[nr][nc] != -1) continue;
            comp[nr][nc] = id;
            st.push_back(make_pair(nr, nc));
        }
    }
    return size;
}

// (3) BFS: same marking rule, a queue instead of a stack (Chapter 3 needs this one)
int fillQueue(int sr, int sc, int id, vector<vector<int>>& comp) {
    queue<pair<int,int>> q;
    comp[sr][sc] = id;
    q.push(make_pair(sr, sc));
    int size = 0;
    while (!q.empty()) {
        int r = q.front().first, c = q.front().second;
        q.pop();
        size++;
        for (int d = 0; d < 4; d++) {
            int nr = r + dr[d], nc = c + dc[d];
            if (!open(nr, nc) || comp[nr][nc] != -1) continue;
            comp[nr][nc] = id;
            q.push(make_pair(nr, nc));
        }
    }
    return size;
}

// (4) scan + fill: label every component, return their sizes
vector<int> labelAll(vector<vector<int>>& comp) {
    comp.assign(n, vector<int>(m, -1));
    vector<int> sizes;
    for (int r = 0; r < n; r++)
        for (int c = 0; c < m; c++)
            if (g[r][c] == '.' && comp[r][c] == -1)
                sizes.push_back(fillStack(r, c, sizes.size(), comp));
    return sizes;
}

int main() {
    n = 4; m = 5;
    g = {"..##.",
         "#.#..",
         ".###.",
         "..#.#"};
    vector<vector<int>> comp;
    vector<int> sizes = labelAll(comp);
    printf("%d components, sizes:", (int)sizes.size());
    for (int s : sizes) printf(" %d", s);                  // 4 components, sizes: 3 4 3 1
    printf("\n");
    for (int r = 0; r < n; r++, printf("\n"))
        for (int c = 0; c < m; c++)
            printf("%c", comp[r][c] < 0 ? '#' : char('A' + comp[r][c]));

    vector<vector<int>> c2(n, vector<int>(m, -1)), c3(n, vector<int>(m, -1));
    dfs(2, 0, 0, c2);                                       // recursive version
    printf("BFS from (0,4): %d cells\n", fillQueue(0, 4, 0, c3));   // 4 cells
}
\end{lstlisting}

\begin{insightbox}[Check the template by hand]
The grid in \texttt{main()} is the one of the scan picture. The labels printed are
\texttt{AA\#\#B / \#A\#BB / C\#\#\#B / CC\#D\#}: A $=\{(0,0),(0,1),(1,1)\}$, B $=\{(0,4),(1,3),(1,4),(2,4)\}$, C $=\{(2,0),(3,0),(3,1)\}$, D $=\{(3,3)\}$ --- sizes 3, 4, 3, 1. Note that $(1,1)$ and $(2,0)$ touch only at a corner, so A and C stay apart.
\end{insightbox}

\textbf{Reading grids.} Rows without spaces (\texttt{\#..\#}) are read with \texttt{cin >> row} into a \texttt{string}; numbers separated by spaces (\texttt{0 1 1}) need \texttt{cin >> a[r][c]} into an \texttt{int} array. For $10^6$ cells add \texttt{ios::sync\_with\_stdio(false); cin.tie(nullptr);} or use \texttt{scanf}.

%------------------------------------------------------------
\subsubsection*{Common mistakes}

\begin{warningbox}[Mistake 1 --- reading the cell before checking the bounds]
\begin{lstlisting}[style=snippet]
if (g[nr][nc] == '.' && nr >= 0 && nr < n && nc >= 0 && nc < m)   // WRONG: g[-1][..] is read first
if (nr >= 0 && nr < n && nc >= 0 && nc < m && g[nr][nc] == '.')   // RIGHT: && stops early
\end{lstlisting}
Out-of-range reads on a \texttt{vector<string>} may ``work'' locally and crash on the judge.
\end{warningbox}

\begin{warningbox}[Mistake 2 --- marking a cell when it is popped instead of when it is pushed]
If \texttt{visited} is set only on \texttt{pop}, a cell next to several marked cells is pushed several times and \emph{processed} several times. On a $2\times2$ island the buggy fill reports area 5, on a $4\times4$ island 25 instead of 16 (we ran it). Mark on push --- every cell enters the container once. (Marking on pop is correct only together with an ``\texttt{if (visited) continue;}'' right after the pop, at up to 4 times the memory.)
\end{warningbox}

\begin{warningbox}[Mistake 3 --- 4 versus 8 neighbours]
\begin{center}
\begin{tikzpicture}
    \begin{scope}
        \gIpIIIgrid{Awww,wBBw,wBww,Cwww}
        \node[font=\scriptsize\bfseries] at (1.05,-3.05) {4-connected: 3 islands};
    \end{scope}
    \begin{scope}[xshift=4.6cm]
        \gIpIIIgrid{Awww,wAAw,wAww,Awww}
        \node[font=\scriptsize\bfseries] at (1.05,-3.05) {8-connected: 1 island};
    \end{scope}
\end{tikzpicture}
\end{center}
Corner contacts join regions only under 8-connectivity. Using the wrong arrays gives a plausible but wrong count --- the samples often contain no diagonal contact, so test one yourself.
\end{warningbox}

\begin{warningbox}[{Mistake 4 --- \texttt{r * n + c} instead of \texttt{r * m + c}}]
The id multiplies by the row \emph{length} (the number of columns). On square grids both are equal, so square samples pass and rectangular tests fail.
\end{warningbox}

\begin{warningbox}[{Mistake 5 --- recursion on a $1000\times1000$ grid}]
See the warning above: the depth can reach $nm$. Use the explicit stack (template 2) or BFS (template 3) as soon as $nm$ exceeds about $10^5$.
\end{warningbox}

\begin{warningbox}[Mistake 6 --- recolouring with the same colour]
In a ``recolour the region'' fill, the visited test is usually ``still has the old colour''. If the new colour equals the old one, nothing ever changes and the recursion never stops. Return immediately when the two colours are equal (Exercise E3.2).
\end{warningbox}

\begin{checklistbox}[``I'm stuck'' checklist for Pattern 3]
\begin{itemize}[leftmargin=*]
    \item[$\square$] What is a vertex: every cell, or only the open / land / same-colour cells?
    \item[$\square$] Which cells are neighbours: 4 sides, 8 directions, knight moves, wrap-around?
    \item[$\square$] Is the question about \emph{regions} (count, size, sum, shape)? Then scan + fill. About \emph{steps}? Then BFS with distances (Chapter 3).
    \item[$\square$] Do I check \emph{inside} before reading the cell, and mark on \emph{push}?
    \item[$\square$] May I overwrite the input as my \texttt{visited} array, or do I need it afterwards?
    \item[$\square$] $nm$ up to $10^6$? Then no recursion, fast input, and \texttt{long long} if I sum values.
    \item[$\square$] Do I need information about a component \emph{later}? Then label every cell with a component id first and store per-id data.
\end{itemize}
\end{checklistbox}
